\section*{How to read this exposition}
\addcontentsline{toc}{section}{How to read this exposition}
The question concerns \emph{distinct vertices}, not the number of paths.
The proof answers a local question about one vertex by introducing sets of
vertices, then sets of ordered pairs, then matrices. Each change of language
has a specific purpose. Write $F(S)=\{y:x\to y\text{ for some }x\in S\}$
and $F^2(S)=F(F(S))$; these are set images with no subtraction.
Keep the following chain in view.
\begin{figure}[ht]
\centering
\begin{tikzpicture}[>=Stealth,box/.style={draw,rounded corners,text width=10.8cm,align=center,inner sep=6pt},node distance=7mm]
\node[box] (a) {A counterexample: every vertex has fewer exact second neighbors};
\node[box,below=of a] (b) {Minimality and blowup: $|U|+|F^2(U)|<2|F(U)|$ for every nonempty proper $U$};
\node[box,below=of b] (c) {Maximize $|P|+|Q|+d$ over compatible pair families containing the diagonal};
\node[box,below=of c] (d) {Complement second images, then prune: obtain a strictly larger objective};
\draw[->] (a)--(b);\draw[->] (b)--(c);\draw[->] (c)--(d);
\end{tikzpicture}
\caption{The dependency chain. The last step uses the pruning lemma, which is proved independently for arbitrary finite relations.}
\end{figure}

\paragraph{Choose a reading route.}
For a five-minute overview, follow the dependency figure, the main theorem,
and the accounting diagram in Section~\ref{sec:incompatibility}.
For a longer first reading, follow the two passes below, accepting pruning
before examining its proof. For a complete verification, use the correctness
ledger and the Lean status document alongside the formal statements. These
routes are adapted from the reading levels in the DeepSeep4.1 rewrite.

\paragraph{Check the ambient set as well as the subsets.}
The main theorem, the graph produced by the reduction, and the no-concavity
proposition all require a nonempty vertex set. A universal claim about
nonempty proper subsets says nothing when the ambient set is empty. The
pruning lemma, in contrast, is valid for an empty ambient set too.

\paragraph{A first pass.}
Read the definitions, the two reduction statements, the pruning statement,
and the extremal proof. On this pass accept the pruning lemma as a tool.
Its unusual cost is the key: deletion is charged partly to points uncovered
\emph{after} pruning. Then return to the linear algebra proving that cost.

\paragraph{A close pass.}
In the reduction, check the exact gain/loss identity for $g(T)$ and the
full-support exception in the blowup. In pruning, check which coordinate
moves forward and which backward, and why all generic rank conditions can
hold at once. In the final proof, check that every fiber is proper and the
diagonal survives. These are the places where an innocent change of
convention would invalidate the argument.

\subsection*{A small example that fixes the conventions}
\begin{figure}[ht]
\centering
\begin{tikzpicture}[>=Stealth,v/.style={circle,draw,minimum size=8mm},every node/.style={font=\small}]
\node[v] (v) at (0,0) {$v$};\node[v] (a) at (2,1) {$a$};
\node[v] (b) at (2,-1) {$b$};\node[v] (c) at (4,0) {$c$};
\draw[->] (v)--(a);\draw[->] (v)--(b);\draw[->] (a)--(b);
\draw[->] (a)--(c);\draw[->] (b)--(c);
\node[align=left,anchor=west] at (5,0) {$N_1^+(v)=\{a,b\}$\\$F^2(\{v\})=\{b,c\}$\\$N_2^+(v)=\{c\}$};
\end{tikzpicture}
\caption{All arcs in this example are drawn. Two paths reach $c$, but it is counted once. The vertex $b$ is reached in two steps and in one step, so it is excluded from the exact second neighborhood. This graph is not a counterexample: $c$ is a sink.}
\end{figure}
For an arbitrary set $S$, the image $F(S)$ may meet $S$. For example,
on a directed triangle and $S$ equal to the whole vertex set, $F(S)=S$.
Never replace a set-image by its exterior neighborhood without explicitly
subtracting $S$. For a singleton in an oriented graph, $v\notin F^2(\{v\})$
because a return in two steps would require opposite arcs. Thus the source's
formula $N_2^+(v)=F^2(\{v\})\setminus F(\{v\})$ is valid.

\Needspace{13\baselineskip}
\paragraph{Notation dictionary.}
\begin{center}
\begin{tabular}{p{2.5cm}p{9cm}}
$D$, then $G$ & Minimal counterexample, then its tournament blowup.\\
$F$, $F^2$ & Union of out-neighbors; its iterate, with no subtraction.\\
$P,Q$ & Compatible left and right families of ordered pairs.\\
$F_1,F_2$ & Advance the first or second coordinate along an arc.\\
$d$ & Points in neither coordinate image of the current pair.\\
$\mathcal R,\mathcal C$ & Prospective families before pruning.\\
$H,Z$ & The two opposite intermediate corners of conflicts.\\
$\Phi$, $\Phi^{p,i}$ & Matrices in the pruning proof; $D$ remains a graph.\\
$L,N,B,A$ & Linear maps in the commuting square used only in pruning.
\end{tabular}
\end{center}
The symbols $I,E,T,Q,C$ in the subset-deficit proof are local to that proof;
in particular its set $Q$ is unrelated to the later pair family $Q$.
